\section{Local capacities and decreasing products}
\label{sec:hessian-smooth-capacity}

The codimension-two example
\cref{hessian:sobolev-two-product-obstruction} shows why Sobolev
capacity cannot control a product of Hessians. Local capacities built
from smooth Hessian measures give the stronger control needed for the
recursive currents. Convergence in these capacities, and hence
quasicontinuity, comes from an energy variation that retains the
derivatives of the Hermitian background.

Work on a coordinate ball $U$ with compact closure in a holomorphic
chart, with the restriction of $\omega$. A smooth local
$m$-subharmonic function means a smooth function $f$ on $U$ that is
$m$-subharmonic relative to $\omega$ in the sense defined before
\cref{hessian:local-unshifted-correction}, that is,
$\ddc f\in\overline{\Gamma_m(\omega)}$ pointwise. For such $f$ put
\[
 \Lambda_s(f)=(\ddc f)^s\wedge\omega^{n-s},\qquad0\leq s\leq m,
\]
where $\Lambda_0(f)=\omega^n$. The local mass and gradient estimates
needed for the capacity use ordinary smooth products.

\begin{lemma}\label{hessian:local-smooth-bounds}
Fix $K\Subset U$ and $M<\infty$. For smooth local $m$-subharmonic
$f_1,\ldots,f_q$ with $|f_i|\leq M$, $0\leq q\leq m$, the mass on $K$
of their mixed trace
$\ddc f_1\wedge\cdots\wedge\ddc f_q\wedge\omega^{n-q}$ is at most a
constant depending only on $K,U,\omega,m,M$. If also $f$ is smooth
local $m$-subharmonic with $|f|\leq M$, then, for $q\leq m-1$ and a
constant $C$ depending on the same data,
\[
 \int_K\mathrm df\wedge\dc f\wedge
       \ddc f_1\wedge\cdots\wedge\ddc f_q
                         \wedge\omega^{n-q-1}\leq C.
\]
\end{lemma}
\begin{proof}
For one potential $f$, the form $\omega+\ddc f$ lies in
$\Gamma_m(\omega)$, and $(\omega+\ddc f)^q\wedge\omega^{n-q}$ expands
into nonnegative terms, one of which is $(\ddc f)^q\wedge\omega^{n-q}$.
The Chern--Levine--Nirenberg bound in the proof of
\cref{hessian:continuous-products}, whose proof uses only the values of
the potential on $U$, therefore bounds the mass on $K$ by $C(1+M)^q$.
For mixed factors, positivity of every term in the expansion of
$(\sum_i\ddc f_i)^q\omega^{n-q}$ bounds the mixed trace by a fixed
multiple of the single trace for $\sum_i f_i$. This proves the first
assertion, including its degree-zero convention.

Replace $f$ by $f+M$, so that $0\leq f\leq2M$ without changing its
gradient or Hessian. The identity
\[
 \ddc(f^2)=2f\ddc f+2\mathrm df\wedge\dc f
\]
also shows that $f^2$ is smooth local $m$-subharmonic. Wedge it with
the remaining $q$ Hessians and $\omega^{n-q-1}$. The first term on
the right then becomes nonnegative, since $f\geq0$ and the degree is
$q+1\leq m$. Hence the gradient term is at most the mixed trace of
$f^2,f_1,\ldots,f_q$, whose mass on $K$ is bounded by the first
assertion with $\max\{M,4M^2\}$ in place of $M$.
\end{proof}

\Needspace{16\baselineskip}
\begin{definition}\label{hessian:smooth-local-capacity}\index{Capacity!smooth local Hessian capacity}
For Borel $E\subset U$ define the smooth capacity
\[
 c_s(E)=\sup\left\{\int_E \Lambda_s(a):
        a\in C^\infty(U),\ -1\leq a\leq0,
        \ddc a\in\overline{\Gamma_m(\omega)}\right\},
\]
for $0\leq s\leq m$. Its outer-open version is
\[
 c_s^*(E)=\inf\{c_s(G):E\subset G\subset U,\ G\text{ open}\}.
\]
A function on a compact $K\Subset U$ is quasicontinuous
for this capacity if, for every $\epsilon>0$, it is continuous on
$K\setminus G$ for some open $G\subset U$ with $c_s(G)<\epsilon$.
\end{definition}

In particular $c_0(E)=\int_E\omega^n$. Each $c_s$ is monotone,
countably subadditive and continuous under increasing unions.
For the last assertion, the supremum over competitors commutes with
the increasing supremum over sets. By
\cref{hessian:local-smooth-bounds}, $c_s$ is finite on compact
subsets of $U$. The outer capacity $c_s^*$ is monotone
and countably subadditive and agrees with $c_s$ on open sets.

The capacities control mixed traces for smooth local $m$-subharmonic functions
$f_1,\ldots,f_q$ with a common bound. Choose a fixed smooth strictly
plurisubharmonic $b$ with $\ddc b\geq\omega$ on $U$. For $q\leq s$,
choose constants $B$ and $L\geq1$ so that
$(f_1+\cdots+f_q+b-B)/L$ lies in $[-1,0]$.
Both constants depend only on the common bound and the fixed data.
Positivity of the smooth binomial expansion gives
\begin{equation}\label{eq:hessian:smooth-capacity-bound}
 (\ddc f_1+\cdots+\ddc f_q)^q\omega^{n-q}
       \leq C \Lambda_s\bigl((\textstyle\sum_i f_i+b-B)/L\bigr).
\end{equation}
Indeed, the expansion of the right side contains the term
\[
 (\ddc f_1+\cdots+\ddc f_q)^q(\ddc b)^{s-q}\omega^{n-s}
\]
with a positive coefficient, and replacing any of its $s-q$ factors
$\ddc b$ by $\omega$ subtracts a nonnegative mixed trace.
For $q=0$ the left side is $\omega^n$, and the same argument applies.
Thus every such mixed trace, and every coefficient of a smooth
product of $q<m$ Hessians, is dominated on Borel sets by $C c_s$.
For the coefficients, this follows from the absolute mixed-form
estimate \cref{hessian:absolute-mixed-forms}.

\Needspace{16\baselineskip}
\begin{lemma}\label{hessian:smooth-torsion}\index{Torsion!local Hessian estimates}
Let $0\leq q\leq m-1$ and let $\alpha$ be a smooth real $(1,1)$-form
in $\overline{\Gamma_m(\omega)}$ pointwise.
Put $T=\alpha^q\omega^{n-q-2}$. For smooth real $f$
and nonnegative Borel weights $a$, supported on a fixed compact set,
\[
 \left|\int a\,\mathrm df\wedge\dc\omega\wedge T\right|
 \leq C\left(\int a\,\mathrm df\wedge\dc f\wedge\omega\wedge T\right)^{1/2}
          \left(\int a\,\omega^2\wedge T\right)^{1/2},
\]
where $C$ depends only on $\omega$ and the compact set.
The corresponding inequality with $\dc f\wedge\mathrm d\omega$
holds as well. For a mixed product of $q$ smooth forms in
$\overline{\Gamma_m(\omega)}$ in place of $\alpha^q$, the left side
is bounded by the sum of these bounds over the finitely many
nonempty subset sums of its factors.
\end{lemma}
\begin{proof}
At one point choose an $\omega$-unitary coframe $e_i$
that diagonalizes $\alpha$, with eigenvalues $\lambda\in\overline{\Gamma_m}$. Put $\epsilon_i=\mathrm i e_i\wedge\bar e_i$
and $\Omega=\epsilon_1\cdots\epsilon_n$, so that
$\omega=\sum_i\epsilon_i$ and $\alpha=\sum_i\lambda_i\epsilon_i$.
Writing $B_{ij}=S_{q;ij}(\lambda)$ and $D_i=S_{q;i}(\lambda)$,
ordinary wedge expansion gives
\[
 \alpha^q\omega^{n-q-2}
 =q!(n-q-2)!\sum_{i<j}B_{ij}
                           \prod_{k\notin\{i,j\}}\epsilon_k.
\]
Only the remaining two-plane contributes to its wedge with
$\mathrm df\wedge\dc\omega$. If $\partial f=\sum_i f_i e_i$,
its absolute density relative to $\Omega$ is bounded by
\[
 C\sum_{i<j}(|f_i|+|f_j|)|B_{ij}|.
\]
The constant bounds the tensor coefficients of $\dc\omega$ in all
unitary frames on the fixed compact set; no derivative of this
pointwise frame is taken.

For $q\geq1$, \cref{hessian:deleted-newton} applies, since
$q\leq m-1\leq n-2$, and gives $|B_{ij}|\leq C(D_iD_j)^{1/2}$; for
$q=0$, $B_{ij}=D_i=1$. Finite Cauchy--Schwarz therefore bounds this sum by
\[
 C\sum_i|f_i|\sqrt{D_i}\sum_j\sqrt{D_j}
 \leq C\left(\sum_i|f_i|^2D_i\right)^{1/2}
          \left(\sum_iD_i\right)^{1/2}.
\]
\Needspace{7\baselineskip}
Here $\sum_iD_i=(n-q)S_q(\lambda)$. Moreover
\[
 \begin{aligned}
 \alpha^q\omega^{n-q}&=q!(n-q)!S_q(\lambda)\Omega,\\
 \mathrm df\wedge\dc f\wedge \alpha^q\omega^{n-q-1}
   &=\kappa q!(n-q-1)!\sum_i|f_i|^2D_i\,\Omega,
 \end{aligned}
\]
where $\kappa=1/(2\pi)$ by the book's $\dc$ normalization.
These two densities are nonnegative. Thus the pointwise torsion
density is bounded by a constant times their geometric mean.
Multiplication by $a$ and integral Cauchy--Schwarz give the assertion,
including the cases where either density vanishes. The other wedge
is conjugate up to a fixed sign, which does not affect its absolute
value.

For the mixed assertion with $q\geq1$, use the exact polarization
identity
\[
 q!\,\alpha_1\wedge\cdots\wedge\alpha_q
   =\sum_{J\subset\{1,\ldots,q\}}
           (-1)^{q-|J|}\left(\sum_{i\in J}\alpha_i\right)^q.
\]
Every nonempty subset sum remains in the cone. Applying the single
power estimate to each term and taking absolute values proves the
mixed assertion. At $q=0$ use the single term $\alpha^0=1$.
\end{proof}

\begin{lemma}\label{hessian:boundary-localization}
Let $f_j$ be uniformly bounded smooth local $m$-subharmonic functions
with $f_j\downarrow f$, and fix a smaller concentric ball $V\Subset U$.
There are uniformly bounded smooth local $m$-subharmonic functions
$v_j\downarrow v$ which equal $f_j$ on a neighborhood of $\overline V$
and equal one fixed strictly plurisubharmonic $\psi$ on one fixed
collar $\{R-\delta_0<|z|<R\}$ of the boundary of $U=\{|z|<R\}$.
Moreover $v=f$ near $\overline V$ and $v=\psi$ on that collar.
\end{lemma}
\begin{proof}
Let $|f_j|\leq M$. Take $\psi=A(|z|^2-R^2)+M+3$, with $A$ so large
that $\psi<-M-2$ on a neighborhood of $\overline V$. Since $\psi=M+3$
on $\partial U$, it exceeds $M+2$ on a sufficiently thin collar
$\{R-\delta_0<|z|<R\}$.

Use the regularized maximum $M_\epsilon$ constructed in
\cref{hessian:regularized-maximum}. Choose $0<\epsilon_j\downarrow0$
with $\epsilon_j<1$, and put $v_j=M_{\epsilon_j}(f_j,\psi)$.
That lemma shows that $v_j$ is smooth local $m$-subharmonic. The
monotonicity of $M_\epsilon$ in its arguments and in $\epsilon$ gives
$v_j\downarrow\max(f,\psi)=v$. Since $f_j-\psi>2>\epsilon_j$ near
$\overline V$ and $\psi-f_j>2>\epsilon_j$ on the collar, the two
asserted regions of exact equality follow. The bounds
\[
 \max(f_j,\psi)\leq v_j\leq\max(f_j,\psi)+\epsilon_j/2
\]
give a common bound. Since every $v_j$ equals $\psi$ on the collar,
all differences $v_k-v_j$ have support in the fixed compact set
$\{|z|\leq R-\delta_0\}$.
\end{proof}

\Needspace{14\baselineskip}
\begin{proposition}\label{hessian:smooth-capacity-convergence}\index{Quasicontinuity!Hessian capacity}
For the sequence $v_j\downarrow v$ of
\cref{hessian:boundary-localization} and $0\leq s\leq m$,
\[
 c_s\{v_k-v>\delta\}\longrightarrow0
                 \quad(k\to\infty)\qquad\text{for every }\delta>0.
\]
Consequently the pointwise limit of a uniformly bounded decreasing
sequence of smooth local $m$-subharmonic functions is quasicontinuous
on every compact subset of $U$, for each $c_s^*$ as well as $c_s$,
$0\leq s\leq m$.
\end{proposition}
\begin{proof}
The assertion for $s=0$ follows by bounded convergence for volume.
Induct on $s$, assuming it for $s-1$, and put $q=s-1$ and
$\theta=\omega^{n-s}$. For $j\geq k$ let $h=v_k-v_j\geq0$.
The support of $h$ lies in one fixed compact set $K_0\Subset U$.
The capacity bound \eqref{eq:hessian:smooth-capacity-bound} and the
induction hypothesis give
\begin{equation}\label{eq:hessian:lower-weight}
 \int h\,\ddc g_1\cdots\ddc g_q\omega^{n-q}
                          \longrightarrow0\quad(k\to\infty),
\end{equation}
uniformly in $j\geq k$ and in the smooth local $m$-subharmonic
$g_1,\ldots,g_q$ with a fixed common bound.
To see this, split at $h=\delta$. On $\{h\leq\delta\}$ the mass
bound of \cref{hessian:local-smooth-bounds} gives $C\delta$. On the
remaining set use $h\leq v_k-v$, the capacity bound with $s-1$ in
place of $s$, and the induction hypothesis. Let $k$ tend to
infinity before $\delta\downarrow0$. For $q=0$ this is the volume
estimate.

Consider the real smooth energy
\[
 E_s(w)=\int_U w(\ddc w)^s\theta,
 \qquad w_t=v_j+t h,\quad0\leq t\leq1.
\]
Each $w_t=(1-t)v_j+tv_k$ is smooth local $m$-subharmonic, with the
common bound of the $v_j$. The energies $E_s(v_j)$ are bounded: on
the fixed collar the integrand is fixed, and on the remaining compact
set \cref{hessian:local-smooth-bounds} gives a uniform mass bound.
Since $h$ vanishes near the boundary,
integration by parts and closedness of $\ddc w_t$ give the exact
identity
\begin{equation}\label{eq:hessian:energy-variation}
 \begin{split}
 \frac{\mathrm d}{\mathrm dt}E_s(w_t)
  &=(s+1)\int h(\ddc w_t)^s\theta\\
  &\quad+s\int h(\ddc w_t)^q
       \bigl(\mathrm dw_t\wedge\dc\theta
           -\dc w_t\wedge\mathrm d\theta+w_t\ddc\theta\bigr).
 \end{split}
\end{equation}
Here
\[
 \begin{aligned}
 \mathrm d\theta&=(n-s)\mathrm d\omega\wedge\omega^{n-s-1},\\
 \ddc\theta&=(n-s)\ddc\omega\wedge\omega^{n-s-1}\\
 &\quad+(n-s)(n-s-1)\mathrm d\omega\wedge\dc\omega
                                           \wedge\omega^{n-s-2},
 \end{aligned}
\]
and the analogous $\dc$ formula holds. A zero scalar coefficient
means that its term is absent. As $k\to\infty$, the two first-order
terms in \eqref{eq:hessian:energy-variation} tend to zero uniformly
in $t$ and $j\geq k$, by \cref{hessian:smooth-torsion}, the gradient
bound of \cref{hessian:local-smooth-bounds}, and
\eqref{eq:hessian:lower-weight}. The last term has the same
property by the absolute-form estimate at degree $q<m$ and
\eqref{eq:hessian:lower-weight}.
More explicitly, the first torsion integral is bounded by
\[
 C\left(\int h\,\mathrm dw_t\wedge\dc w_t
                     \wedge(\ddc w_t)^q\omega^{n-q-1}\right)^{1/2}
   \left(\int h(\ddc w_t)^q\omega^{n-q}\right)^{1/2}.
\]
The first factor is bounded since $h$ has a common bound and
\cref{hessian:local-smooth-bounds} applies; the second tends to
zero by \eqref{eq:hessian:lower-weight}. Conjugation handles the
other first-order term.
Integrating \eqref{eq:hessian:energy-variation} over $0\leq t\leq1$
therefore gives, for $j\geq k$,
\[
 E_s(v_k)-E_s(v_j)=(s+1)\Delta_{kj}+r_{kj},
 \quad
 \Delta_{kj}=\int_0^1\int h\Lambda_s(w_t)\,\mathrm dt\geq0,
 \quad |r_{kj}|\leq e_k,
\]
where $e_k\to0$ does not depend on $j$.
It follows that $E_s(v_j)\leq E_s(v_k)+e_k$ for $j\geq k$, so
$\limsup_jE_s(v_j)\leq E_s(v_k)+e_k$ for every $k$. Letting $k$ run
through a subsequence along which $E_s(v_k)$ tends to its lower
limit shows that the bounded sequence $E_s(v_j)$ converges. Hence
$\Delta_{kj}\to0$ as $k\to\infty$, uniformly for $j\geq k$.
Positivity of the binomial expansion
$\Lambda_s(w_t)\geq(1-t)^s \Lambda_s(v_j)$ then gives
\begin{equation}\label{eq:hessian:self-weight}
 \int h\Lambda_s(v_j)\longrightarrow0
             \quad(k\to\infty),\ \text{uniformly for }j\geq k.
\end{equation}

It remains to replace each Hessian factor by a capacity competitor.
For any smooth local $m$-subharmonic $a$ with $-1\leq a\leq0$, put
\[
 \begin{aligned}
 I_\ell&=\int h(\ddc v_j)^\ell(\ddc a)^{s-\ell}\theta,
   &&0\leq\ell\leq s,\\
 \eta_\ell&=(\ddc v_j)^\ell(\ddc a)^{s-\ell-1},
   &&0\leq\ell<s.
 \end{aligned}
\]
Stokes' theorem for the compactly supported form
$h\,\dc a\wedge\eta_\ell\theta$, with every metric derivative
kept, gives
\[
 I_\ell=-\int\mathrm dh\wedge\dc a\wedge\eta_\ell\theta
          +\int h\,\dc a\wedge\eta_\ell\mathrm d\theta.
\]
The second term is $o(1)$ as $k\to\infty$, uniformly in $a$ and
$j\geq k$. Indeed, apply the mixed form of
\cref{hessian:smooth-torsion} with $f=a$ and weight $h$; each subset
sum has a uniform gradient bound by \cref{hessian:local-smooth-bounds},
and its lower trace is controlled by \eqref{eq:hessian:lower-weight}.

The form $\eta_\ell\theta$ is positive of bidegree $(n-1,n-1)$:
a complementary nonnegative rank-one form completes a polarized
product of $s\leq m$ cone forms. Its associated Hermitian
gradient pairing is nonnegative, so Cauchy--Schwarz and the uniform
gradient bound for $a$ give
\[
 I_\ell\leq C F_\ell(h)^{1/2}+o(1),
 \qquad F_\ell(h)=\int\mathrm dh\wedge\dc h\wedge\eta_\ell\theta.
\]
Stokes' theorem for $h\,\dc h\wedge\eta_\ell\theta$ gives
\[
 F_\ell(h)= I_{\ell+1}
       -\int h\ddc v_k\wedge\eta_\ell\theta
       +\int h\,\dc h\wedge\eta_\ell\mathrm d\theta
                    \leq I_{\ell+1}+o(1).
\]
The dropped term is nonnegative. For the last term the same
polarized torsion estimate applies, with $f=h$ and weight $h$. Its
first gradient factor is uniformly bounded because $0\leq h\leq C$
and $h=v_k-v_j$: for every positive single-power gradient pairing,
the gradient energy of the difference is at most twice the sum of
those of $v_k$ and $v_j$. Its second factor tends to zero by
\eqref{eq:hessian:lower-weight}.
Starting with \eqref{eq:hessian:self-weight} at $I_s$, finite
backward induction gives $I_0\to0$ uniformly in $a$ and $j\geq k$.

Since $I_0=\int h\Lambda_s(a)$, Chebyshev's inequality now yields
$c_s\{v_k-v_j>\delta\}\leq\delta^{-1}\sup_a I_0$.
These sets increase to $\{v_k-v>\delta\}$ as $j\to\infty$.
Continuity of $c_s$ under increasing unions proves the induction.

For quasicontinuity, fix $0\leq s\leq m$ and $\epsilon>0$, and
choose $k_N$ so that the sets
$G_N=\{v_{k_N}-v>2^{-N}\}$ satisfy
$c_s(G_N)<\epsilon2^{-N}$, $N\geq1$.
They are open because $v$ is upper semicontinuous and $v_{k_N}$
is smooth. On $U\setminus\bigcup_NG_N$ these smooth functions
converge uniformly to the pointwise limit $v$. Hence $v$ is
continuous there. The union has capacity less than $\epsilon$.
Every compact subset of $U$ lies in a smaller concentric ball $V$,
and $v=f$ near $\overline V$ in \cref{hessian:boundary-localization};
this transfers the conclusion to $f$. Open sets have the same outer
capacity, which also proves the assertion for $c_s^*$.
\end{proof}

\Needspace{16\baselineskip}
\begin{lemma}\label{hessian:rough-coefficient-capacity}
Every bounded admissible potential on $X$ is locally quasicontinuous
for each smooth capacity $c_s^*$, $0\leq s\leq m$. For $0\leq p<m$,
each coefficient measure $\nu$ of the ordered $D_p$ satisfies
\[
 |\nu|(E)\leq C c_{m-1}^*(E),\qquad E\subset K\Subset U,
\]
where $C$ depends only on $K$, $U$, $\omega$ and a common bound for
the potentials. For $p=0$ the assertion concerns the volume measure,
the coefficient of $D_0=1$.
\end{lemma}
\begin{proof}
Let $u$ be bounded admissible. The global decreasing approximation
of \cref{hessian:bounded-decreasing-approximation} and
\cref{hessian:local-unshifted-correction} provide uniformly bounded
smooth local $m$-subharmonic functions decreasing to $u+\rho$ on $U$.
By \cref{hessian:smooth-capacity-convergence}, this pointwise limit,
formed with the canonical values of $u$, is quasicontinuous.
Subtracting the smooth $\rho$ preserves quasicontinuity.

For smooth admissible potentials $u_i$ write
$\ddc u_i=\ddc(u_i+\rho)-\ddc\rho$ and expand the product.
Each term has $q\leq p<m$ local cone factors and fixed smooth
complementary factors. The absolute-form estimate and
\eqref{eq:hessian:smooth-capacity-bound}, with $s=m-1$, bound
each coefficient variation by $C c_{m-1}$ on Borel sets.
The constants depend only on the common bound and local data.

Take the iterated limits of
\cref{hessian:bounded-recursive-currents} on a slightly larger
compact neighborhood of $K$. Lower semicontinuity of signed
measure variation on open sets preserves these estimates at each
limit. For any open $G\supset E$, intersect with a fixed open
neighborhood of $K$ having compact closure in $U$; this gives
$|\nu|(E)\leq Cc_{m-1}(G)$. Taking the infimum proves the outer
estimate.
\end{proof}

\begin{lemma}\label{hessian:weighted-measure-convergence}
Fix $0\leq s\leq m$ and the outer smooth capacity $c_s^*$ on $U$.
Let signed Radon measures $\nu_j\to\nu$ weakly have uniformly bounded
local variations, with all variations dominated on compact subsets
by $C c_s^*$. Suppose, for $1\leq i\leq\ell$, that bounded Borel
functions $b_i^j\downarrow b_i$ have one common bound, and that every
$b_i^j$ and every limit $b_i$ is locally quasicontinuous for this
capacity. For any continuous real $F$ on a compact box containing
their ranges,
\[
 F(b_1^j,\ldots,b_\ell^j)\nu_j
       \longrightarrow F(b_1,\ldots,b_\ell)\nu
\]
weakly as signed Radon measures.
\end{lemma}
\begin{proof}
Fix $\phi\in C_c(U)$, a compact $K\Subset U$ containing its
support, and $\epsilon>0$. For each of the countably many functions
$b_i^j$ and $b_i$, choose an exceptional open set for
quasicontinuity on $K$, with capacities summing to less than
$\epsilon$. Their union $G$ is open, $c_s(G)<\epsilon$, and every
function is continuous on the compact set $A=K\setminus G$.
Dini's theorem gives $b_i^j\to b_i$ uniformly on $A$ for each $i$.
Uniform continuity of $F$ on the compact box therefore gives
$f_j\to f$ uniformly on $A$, where
$f_j=F(b_1^j,\ldots,b_\ell^j)$ and $f=F(b_1,\ldots,b_\ell)$.
Put $M=\|F\|_\infty$. Extend $f|_A$ to a continuous $g$ on $U$
with $|g|\leq M$; if $A$ is empty, take $g=0$.
Then
\[
\begin{aligned}
 \left|\int\phi f_j\,\mathrm d\nu_j-
                 \int\phi f\,\mathrm d\nu\right|
 &\leq \|\phi\|_\infty\|f_j-f\|_A|\nu_j|(K)
       +\left|\int\phi g\,\mathrm d(\nu_j-\nu)\right|\\
 &\quad+2M\|\phi\|_\infty
            \bigl(|\nu_j|(K\cap G)+|\nu|(K\cap G)\bigr).
\end{aligned}
\]
For fixed $G$, the first two terms tend to zero as $j\to\infty$:
the first by uniform convergence on $A$, the second by weak
convergence applied to $\phi g\in C_c(U)$. Since
$|\nu_j|(K\cap G)\leq Cc_s(G)$, and similarly for $\nu$, the last
term is at most $4MC\|\phi\|_\infty\epsilon$. Let $\epsilon$ tend
to zero.
\end{proof}

\Needspace{20\baselineskip}
\begin{theorem}\label{hessian:joint-decreasing-products}\label{hessian:bounded-two-factor-convergence}\label{hessian:bounded-linear-weighted-convergence}\index{Hessian measure!simultaneous decreasing convergence}\index{Hessian measure!two bounded factors}\index{Hessian measure!weighted two-factor convergence}
For uniformly bounded admissible potentials $u_i^j\downarrow u_i$
on $X$, the currents $D_p(u_1^j,\ldots,u_p^j)$ converge to
$D_p(u_1,\ldots,u_p)$ in distributions for $p\leq m$, and weakly
with measure coefficients for $p<m$. These currents are symmetric
in their potentials. For $p<m$ the coefficient convergence also
holds after multiplication by a uniformly bounded decreasing
admissible weight. The positive shifted currents $Q_p$ converge
weakly for $p\leq m$; $Q_m$ converges as a positive Radon measure.
\end{theorem}
\begin{proof}
Work on a compact coordinate test set, and induct on $p$; the case
$p=0$ is trivial since $D_0=1$. Suppose the result holds for
$D_{p-1}$, which has measure coefficients since $p-1<m$.
\Cref{hessian:rough-coefficient-capacity,hessian:weighted-measure-convergence},
applied coefficientwise with $s=m-1$ and weight $u_p^j$, give weak
convergence of $u_p^jD_{p-1}(u_1^j,\ldots,u_{p-1}^j)$ to
$u_pD_{p-1}(u_1,\ldots,u_{p-1})$. Applying $\ddc$ proves
distributional convergence at degree $p$. For $p<m$, the uniform
coefficient bounds in \cref{hessian:bounded-recursive-currents}
upgrade this to weak convergence of measures, by smooth
approximation of compactly supported continuous tests.
\Cref{hessian:weighted-measure-convergence}, now with a decreasing
weight $a_j\downarrow a$, gives the additional weighted conclusion.

For each fixed list choose simultaneous smooth decreasing global
approximations. Their smooth raw Hessian wedges are symmetric.
The convergence just proved identifies their limit with $D_p$
in any ordering. This proves symmetry without exchanging the
iterated limits of \eqref{eq:hessian-bounded-iterated-limit}.

The subset expansion for $R_p$ uses fixed smooth powers of
$\omega$, so its distributional convergence and symmetry follow
from those of $D_I$. Its wedge $Q_p$ is positive with uniform
local masses by \cref{hessian:bounded-shifted-positive}.
Distributional convergence and these mass bounds imply weak
convergence as an order-zero positive current, including the
Radon measure conclusion at $p=m$. A partition of unity on a
finite coordinate cover proves the global assertions.
\end{proof}

The coefficient bound for the raw currents $D_p$, $p<m$, uses only
$c_{m-1}$. For the positive shifted products, $c_m$ controls the
traces in every degree $p\leq m$, including the top degree, where
$D_m$ itself has no coefficient bound. This allows a bounded
decreasing weight in every Hessian degree.

\Needspace{15\baselineskip}
\begin{lemma}\label{hessian:positive-capacity-domination}
On $K\Subset U$, the positive shifted currents
$Q_p(u_1,\ldots,u_p)$, $1\leq p\leq m$, for globally bounded
admissible potentials have trace and coefficient variations bounded
on Borel $E\subset K$ by $C c_m^*(E)$. The constants are uniform
when all potentials have a common bound.
\end{lemma}
\begin{proof}
First take smooth potentials, put $\sigma_i=\omega+\ddc u_i$, and
let $\rho$ be the fixed local correction of
\cref{hessian:local-unshifted-correction}, with
$\tau=\ddc\rho\geq2\omega$. Then
$\beta_i=\ddc(u_i+\rho)=\sigma_i+(\tau-\omega)$ lies in
$\overline{\Gamma_m(\omega)}$. Since $\tau-\omega\geq0$, mixed
positivity gives
\[
 0\leq\int_E Q_p\wedge\omega^{m-p}
    \leq\int_E\beta_1\cdots\beta_p\omega^{n-p}
    \leq C c_m(E),\qquad E\subset K.
\]
Here $Q_p\wedge\omega^{m-p}=\sigma_1\cdots\sigma_p\omega^{n-p}$.
The last bound is the domination of mixed traces by $c_m$ that
follows \eqref{eq:hessian:smooth-capacity-bound}, applied with $s=m$
to the corrected potentials $u_i+\rho$, which are bounded on $U$.
Positivity of $Q_p$ bounds each of its
coefficient variations on the compact coordinate set by a fixed
multiple of its trace measure.

For bounded potentials choose uniformly bounded smooth decreasing
approximations by \cref{hessian:bounded-decreasing-approximation}. The joint
convergence in \cref{hessian:joint-decreasing-products} identifies
their positive-current limit with the previously constructed $Q_p$.
On a compact neighborhood of $K$ inside $U$, lower semicontinuity
on open sets preserves trace and coefficient variation estimates.
For an arbitrary open $G\supset E$ intersect with a fixed relatively
compact neighborhood of $K$ before taking this limit; its capacity
is at most $c_m(G)$. Therefore every such trace or variation measure
$\mu$ satisfies $\mu(E)\leq Cc_m(G)$. Infimize over $G$ to obtain
the asserted outer-open bound. The constants of the smooth estimates
are independent of the approximation index.
\end{proof}

\Needspace{15\baselineskip}
\begin{theorem}\label{hessian:weighted-positive-products}\index{Hessian measure!weighted positive products}
Suppose $a_j\downarrow a$ and $u_i^j\downarrow u_i$ are bounded
admissible potentials on $X$ with common bounds. Then,
for $1\leq p\leq m$,
\[
 a_j Q_p(u_1^j,\ldots,u_p^j)
        \longrightarrow a Q_p(u_1,\ldots,u_p)
\]
weakly as currents with measure coefficients. At $p=m$ this is weak
convergence of weighted Radon measures.
\end{theorem}
\begin{proof}
Work on a compact coordinate test set. By
\cref{hessian:joint-decreasing-products}, every coefficient measure
of the positive shifted currents converges weakly and has a uniform
local variation bound. \Cref{hessian:positive-capacity-domination}
strengthens this to domination by $C c_m^*$ on Borel sets.
\Cref{hessian:rough-coefficient-capacity} makes the weights
$a_j$ and $a$ quasicontinuous for this capacity. Apply
\cref{hessian:weighted-measure-convergence} with $s=m$ to each
coefficient, and then use a finite partition of unity. At $p=m$,
$Q_m$ is a positive Radon measure, so this also gives the stated
weighted measure convergence.
\end{proof}
